\documentclass[11pt,a4paper,landscape]{article}

\usepackage{fontspec}
\usepackage[margin=14mm]{geometry}
\usepackage{graphicx}
\usepackage{caption}
\usepackage{xcolor}
\usepackage{hyperref}
\usepackage{enumitem}
\usepackage{float}

\defaultfontfeatures{Ligatures=TeX,Scale=MatchLowercase}
\IfFontExistsTF{Noto Sans CJK KR}{
  \setmainfont{Noto Sans CJK KR}[
    ItalicFeatures={FakeSlant=0.18},
    BoldItalicFeatures={FakeSlant=0.18}
  ]
  \setsansfont{Noto Sans CJK KR}[
    ItalicFeatures={FakeSlant=0.18},
    BoldItalicFeatures={FakeSlant=0.18}
  ]
}{
  \IfFontExistsTF{Apple SD Gothic Neo}{
    \setmainfont{Apple SD Gothic Neo}[
      ItalicFeatures={FakeSlant=0.18},
      BoldItalicFeatures={FakeSlant=0.18}
    ]
    \setsansfont{Apple SD Gothic Neo}[
      ItalicFeatures={FakeSlant=0.18},
      BoldItalicFeatures={FakeSlant=0.18}
    ]
  }{
    \setmainfont{NanumGothic}[
      ItalicFeatures={FakeSlant=0.18},
      BoldItalicFeatures={FakeSlant=0.18}
    ]
    \setsansfont{NanumGothic}[
      ItalicFeatures={FakeSlant=0.18},
      BoldItalicFeatures={FakeSlant=0.18}
    ]
  }
}
\XeTeXlinebreaklocale "ko"
\XeTeXlinebreakskip = 0pt plus 1pt

\hypersetup{
  colorlinks=true,
  linkcolor=blue!50!black,
  urlcolor=blue!50!black,
  pdftitle={UUV MuJoCo v2.2 current visuals},
  pdfauthor={OpenAI Codex}
}

\setlength{\parskip}{0.4em}
\setlength{\parindent}{0pt}
\captionsetup{
  font=small,
  labelfont=bf,
  justification=raggedright,
  singlelinecheck=false
}

\title{UUV MuJoCo v2.2 current 시각 자료집\\
\large 그림 위주 정리본}
\author{}
\date{2026-04-07}

\begin{document}
\maketitle

이 문서는 발표나 회의에서 그림만 빠르게 넘겨보기 위한 별도 자료다. 자세한 설명과 수치 해석은 내용 정리본 PDF를 기준으로 읽는다.

\clearpage
\begin{figure}[H]
  \centering
  \includegraphics[width=0.88\linewidth]{figures_v22_latest/uuv_v22_latest_current_engine_block_diagram.png}
  \caption{current 물리엔진 블록 다이어그램. MuJoCo built-in ellipsoid fluid와 Python custom thrust / hydrostatics의 경계를 한 장에 정리했다.}
\end{figure}

\clearpage
\begin{figure}[H]
  \centering
  \includegraphics[width=0.94\linewidth]{figures_v22_latest/uuv_v22_latest_thruster_model_detail.png}
  \caption{current thrust generation pipeline. 입력 정규화, 1차 지연, deadzone / limit, force map, 축별 응답 보정까지를 순서대로 보여준다.}
\end{figure}

\clearpage
\begin{figure}[H]
  \centering
  \includegraphics[width=0.90\linewidth]{figures_v22_latest/uuv_v22_latest_thruster_performance_curve.png}
  \caption{16V T200 performance reference graph. 좌측은 전체 force curve, 우측은 1500\,us 주변 0-force plateau 축소 구간을 확대해 보여준다.}
\end{figure}

\clearpage
\begin{figure}[H]
  \centering
  \includegraphics[width=0.88\linewidth]{figures_v22_latest/uuv_v22_latest_buoyancy_layout_compare.png}
  \caption{single-point center와 current 4-point buoyancy layout 비교. 같은 총 부력을 어디에 어떻게 분포시키는지의 차이를 시각적으로 보여준다.}
\end{figure}

\clearpage
\begin{figure}[H]
  \centering
  \includegraphics[width=0.94\linewidth]{figures_v22_latest/uuv_v22_latest_buoyancy_variant_responses.png}
  \caption{같은 current scene / thrust model / scripted harness에서 buoyancy 구성만 바꾼 응답 비교. 전진, pitch, heave, attitude release를 한 장에 묶었다.}
\end{figure}

\clearpage
\begin{figure}[H]
  \centering
  \includegraphics[width=0.94\linewidth]{figures_v22_latest/uuv_v22_latest_buoyancy_variant_metrics.png}
  \caption{latest actual reference 대비 key metric과 상대 오차 비교. peak surge, peak pitch / roll, heave depth delta, mean gap을 요약했다.}
\end{figure}

\end{document}
